\src{177}
fixé de \fquote{dimension} $\symup{\leq q}$, alors le foncteur contravariant
$\symup{E(M)=\operatorname{Ext}_{A}^{p}(M,A)}$ sur la catégorie des $\symup{A/J}$-modules est
exact à gauche ; de plus il transforme sommes directes en produits directs,
d'où résulte facilement que
$\symup{E(M)=\operatorname{Hom}_{A}\bigl(M,E(A/J)\bigr)}$.
\footnote{Le texte imprimé porte $\symup{E(A)}$ ; dans la catégorie considérée, l'objet
représentant est $\symup{A/J}$, et l'égalité requise pour (3.7) est $\symup{E(A/J)}$.} Enfin, la
dernière assertion de la proposition 5 est plus subtile, et résulte d'une
caractérisation intrinsèque des $\symup{E^{i}(F)}$ à l'aide d'un théorème de dualité
locale qui ne peut être donné ici.

\result{COROLLAIRE}{} \uline{Désignons par} $\symup{\underline{\omega}_Y^{q}}$ \uline{le faisceau}
$\symup{\Extsh_{\Osh_X}^{p}(\Osh_Y,\underline{\Omega}_X^{n})}$. \uline{Alors on a un isomorphisme
fonctoriel pour les faisceaux algébriques cohérents} $\symup{F}$ \uline{sur} $\symup{Y}$ :
\begin{equation}\tag{3.8}
  \symup{\Extsh_{\Osh_X}^{p}(F,\underline{\Omega}_X^{n})
  =\Homsh_{\Osh_X}(F,\underline{\omega}_Y^{q})}
\end{equation}

\fsection{\normalfont4}{\normalfont\uline{Classe de cohomologie associée à une sous-variété}.}

Dans toute la suite, $\symup{X}$ désigne un ensemble algébrique de dimension $\symup{n}$, défini
sur un corps $\symup{\underline{k}}$, que nous supposerons algébriquement clos pour simplifier. Sauf
au n\textsuperscript{o} 6, $\symup{X}$ est supposée non singulière. On désigne par
$\symup{\Osh_X}$ le faisceau des anneaux locaux sur $\symup{X}$, par
$\symup{\underline{\Omega}_X^{*}=\bigoplus\limits_p\underline{\Omega}_X^{p}}$ le faisceau des germes de formes différentielles sur $\symup{X}$.\footnote{Le texte
imprimé emploie un signe de réunion ; la structure graduée est la somme directe
des faisceaux $\symup{\underline{\Omega}_X^p}$.} Si $\symup{Y}$ est une partie fermée de $\symup{X}$, on identifie
les faisceaux algébriques cohérents sur $\symup{Y}$ à des faisceaux cohérents sur $\symup{X}$
nuls hors de $\symup{Y}$ ; il en est en particulier ainsi de $\symup{\Osh_Y}$ et $\symup{\underline{\Omega}_Y}$.

\result{LEMME}{1} \uline{Soient} $\symup{F}$ \uline{un faisceau algébrique cohérent sur} $\symup{X}$ \uline{dont le
support est de dimension} $\symup{\leq n-p}$, $\symup{L}$ \uline{un faisceau algébrique cohérent
localement libre sur} $\symup{X}$. \uline{Alors} $\symup{\operatorname{Ext}_{\Osh_X}^{i}(X;F,L)}$ \uline{est nul
pour} $\symup{i<p}$, \uline{et on a un isomorphisme canonique}
\begin{equation}\tag{4.1}
  \symup{\operatorname{Ext}_{\Osh}^{p}(X;F,L)
  =H^{0}\bigl(X,\Extsh_{\Osh}^{p}(F,L)\bigr)}
\end{equation}
\uline{Si} $\symup{F}$ \uline{est un faisceau algébrique cohérent sur une partie fermée} $\symup{Y}$ \uline{de} $\symup{X}$ \uline{de
dimension} $\symup{\leq n-p}$, \uline{on a un isomorphisme canonique}
\begin{equation}\tag{4.1 bis}
  \symup{\Extsh_{\Osh_X}^{p}(F,L)
  =\Homsh_{\Osh_X}\!\left(
    F\otimes L'\otimes\underline{\Omega}_X^{n},\underline{\omega}_Y^{n-p}
  \right)}
\end{equation}
\uline{où} $\symup{\underline{\omega}_Y^{n-p}}$ \uline{est le faisceau sur} $\symup{Y}$ \uline{défini dans le corollaire à la
proposition 5, (qui s'identifie à} $\symup{\underline{\Omega}_Y^{n-p}}$ \uline{si} $\symup{Y}$ \uline{est non singulière}).

La formule (4.1) est une conséquence immédiate de la suite spectrale de la
proposition 1, et de la proposition 5 ; en vertu de la formule (3.8) on peut
écrire

